\begin{explanationbox}
	\textbf{\textcolor{CExplanation}{一句话直觉}}
	补集是“反向选择”，补两次就是“否定之否定”，自然回到原集合。

	\medskip
	\textbf{\textcolor{CExplanation}{核心拆解}}
	\begin{description}[style=nextline, leftmargin=2.8em, labelsep=0.5em, itemsep=0.45em]
		\item[\textbf{要点一（补集定义）：}] 在全集中，不属于 $A$ 的元素构成 $A$ 的补集。记号 $\complement_U A$ 就是“$U$ 中挖掉 $A$ 剩下的部分”。
		\item[\textbf{要点二（双补还原）：}] 第一次补得到“圈外”，第二次补又把“圈外”变成了“圈内”，所以两次补运算相当于恒等变换。
		\item[\textbf{要点三（边界特例）：}] 全集的补是没有元素（空集），空集的补是全集，这是补集定义的自然产物。
	\end{description}

	\medskip
	\textbf{\textcolor{CExplanation}{几何本质（韦恩图解）}}
	画一个大长方形代表全集 $U$，内部画一个圆代表 $A$。$\complement_U A$ 是圆外阴影区域；再对这个阴影区域取补，即 $U$ 中除阴影外的部分，恰好是原来的圆 $A$。韦恩图能一眼看出两次取补回到原地。

	\medskip
	\textbf{\textcolor{CExplanation}{代数意义（对合运算）}}
	补运算 $\complement_U$ 作用在集合上，两次作用等于恒等映射：$\complement_U \circ \complement_U = \text{id}$。这类运算称为“对合”，类似于相反数（两次取反回到自身）、逻辑非（非的非等于原命题）。理解这一点，有助于将集合运算与逻辑运算联系起来。

	\medskip
	\textbf{\textcolor{CExplanation}{考点价值（化简与求值）}}
	\begin{itemize}[leftmargin=*, itemsep=0.5em, topsep=0.4em]
		\item \textbf{考法一（直接求双补）：} 已知具体集合，一眼得出 $\complement_U(\complement_U A)=A$，不必逐层计算。
		\item \textbf{考法二（化简表达式）：} 遇到含多层补集的式子，可立即化简，降低出错概率。
	\end{itemize}

	\medskip
	\textbf{\textcolor{CExplanation}{顿悟点}}
	\textbf{把补集看成“非”，双重补集就是“非的非”，逻辑上等于原命题。用这个观点可以瞬间记住公式。}

	\medskip
	\textbf{\textcolor{CExplanation}{使用场景}}
	\begin{itemize}[leftmargin=*, itemsep=0.5em, topsep=0.4em]
		\item \textbf{场景一（已知补求原）：} 题目给出 $\complement_U A$，要求 $A$，直接答 $\complement_U(\complement_U A)$。
		\item \textbf{场景二（化简多重补）：} 在复杂集合关系证明中，遇到 $\complement_U(\complement_U(\cdots))$ 时，直接去掉两层补号。
	\end{itemize}
\end{explanationbox}
